% Options for packages loaded elsewhere
\PassOptionsToPackage{unicode}{hyperref}
\PassOptionsToPackage{hyphens}{url}
%
\documentclass[
  10pt,
]{article}
\usepackage{amsmath,amssymb}
\usepackage{iftex}
\ifPDFTeX
  \usepackage[T1]{fontenc}
  \usepackage[utf8]{inputenc}
  \usepackage{textcomp} % provide euro and other symbols
\else % if luatex or xetex
  \usepackage{unicode-math} % this also loads fontspec
  \defaultfontfeatures{Scale=MatchLowercase}
  \defaultfontfeatures[\rmfamily]{Ligatures=TeX,Scale=1}
\fi
\usepackage{lmodern}
\ifPDFTeX\else
  % xetex/luatex font selection
\fi
% Use upquote if available, for straight quotes in verbatim environments
\IfFileExists{upquote.sty}{\usepackage{upquote}}{}
\IfFileExists{microtype.sty}{% use microtype if available
  \usepackage[]{microtype}
  \UseMicrotypeSet[protrusion]{basicmath} % disable protrusion for tt fonts
}{}
\makeatletter
\@ifundefined{KOMAClassName}{% if non-KOMA class
  \IfFileExists{parskip.sty}{%
    \usepackage{parskip}
  }{% else
    \setlength{\parindent}{0pt}
    \setlength{\parskip}{6pt plus 2pt minus 1pt}}
}{% if KOMA class
  \KOMAoptions{parskip=half}}
\makeatother
\usepackage{xcolor}
\usepackage[margin=0.85in]{geometry}
\setlength{\emergencystretch}{3em} % prevent overfull lines
\providecommand{\tightlist}{%
  \setlength{\itemsep}{0pt}\setlength{\parskip}{0pt}}
\setcounter{secnumdepth}{-\maxdimen} % remove section numbering
\usepackage{amsmath,amssymb,amsthm}
\usepackage{longtable}
\ifLuaTeX
  \usepackage{selnolig}  % disable illegal ligatures
\fi
\usepackage{bookmark}
\IfFileExists{xurl.sty}{\usepackage{xurl}}{} % add URL line breaks if available
\urlstyle{same}
\hypersetup{
  pdftitle={Actual metric frames and canonical traces across nonflat fibers},
  pdfauthor={Primary-source reconstruction; unverified mathematical draft},
  hidelinks,
  pdfcreator={LaTeX via pandoc}}

\title{Actual metric frames and canonical traces across nonflat fibers}
\author{Primary-source reconstruction; unverified mathematical draft}
\date{5 October 2026}

\begin{document}
\maketitle

\subsection{1. Disposition, preserved target, and new
results}\label{disposition-preserved-target-and-new-results}

The unchanged original P16.1 is not solved. Candidate=false. The
original atomic problem is:

\begin{quote}
P16.1 (T5): Finite-energy equations on singular or unibranch varieties
Extend the finite-energy non-Archimedean Monge--Ampère operator and its
existence-and- uniqueness theorem to singular or unibranch varieties,
replacing each use of smoothness in compactness, envelopes, energy
differentiation and comparison. State all retained polarization and field
hypotheses.
\end{quote}

Its completion criterion is preserved verbatim:

\begin{quote}
Establish the genuine finite-energy operator, exact measure mass and
energy condition, existence, uniqueness modulo constants, and every
compactness/envelope/differentiation/comparison replacement for the
singular/geometrically unibranch projective domain with explicitly
retained field and polarization hypotheses. Verify the new envelope
source and the actual general-variety global theory. No change to the
original mathematical target is permitted.
\end{quote}

The complete original problem, family, task object, and assignment are
preserved in \texttt{input.json}, its byte-identical
\texttt{input-original.json}, \texttt{task-original.json}, and
\texttt{assignment-original.txt}. The supplied round72 argument, its
actual four notes, round70, and root80/root81 have been read as
mathematical data, and retained in \texttt{embedded-data/}. Their
printed paths were not accessed. No result from these unreviewed drafts
is silently assumed below.

The equation domain under investigation retains a reduced irreducible
geometrically unibranch projective variety over an arbitrary
algebraically closed trivially valued characteristic-zero field, an
ample line bundle, and arbitrary-support finite-dual-energy measures.
The original problem asks for all replacements in this domain. The
following new partial results do not narrow that target.

\begin{enumerate}
\def\labelenumi{\arabic{enumi}.}
\tightlist
\item
  For a precisely defined continuous rational contact metric that is
  actually glued on an ample algebraic line bundle, we construct a
  strict analytic G-cover, norm-one frames, the original bundle
  identification, and genuine same-bundle FS approximation after exact
  Gauss descent. We prove curve restriction directly from harmonic
  supports. We then construct bounded local products and bounded
  weighted limits on shared genuine collars, independently of supplied
  sequences. Compact IBP and local mass/comparison/domination still
  precede weak globalization.
\item
  A generic determinant on a finite integral algebra over a normal base
  gives a canonical positive integer weight at \textbf{every} analytic
  fiber point, including nonflat fibers. The weights sum to the generic
  degree. Every actual affine FS test has an exact finite determinant
  formula on the whole base. This algebraic result needs neither weak
  regularization nor a line-slicing definition of weights.
\item
  The resulting trace is continuous on the whole base for continuous
  inputs and obeys a degree bound. For continuous AFS-regularizable
  inputs this follows from locally uniform genuine AFS approximants. We
  also prove radial identification of the trace of a weak input,
  \textbf{provided} its open trace already has an actual
  AFS-regularizable extension. No affine regularization theorem over a
  completed residue field is used. Existence of that extension still
  needs Gate B.
\end{enumerate}

These are mathematical proof drafts, rather than execution certificates
or claims of independent review.

\subsection{2. Exact local hypotheses for the metric
result}\label{exact-local-hypotheses-for-the-metric-result}

Write a metric weight as \(b=-\log\|e\|\) in a frame \(e\). A
\emph{unit-chart rational contact weight} on an analytic open \(U\)
means the following data locally at every point:

\begin{itemize}
\tightlist
\item
  finitely many analytic units \(f_1,\ldots,f_r\) on a neighborhood
  \(U'\);
\item
  a continuous function \(F\) on their actual logarithmic image such
  that \(b=F(\log|f_1|,\ldots,\log|f_r|)\);
\item
  \(F\) is a continuous selection from finitely many affine expressions
  with rational slopes and rational constants;
\item
  at each image point \(t\), an affine expression \(\ell_t+C_t\),
  allowed real slopes, is below \(F\) on a neighborhood of \(t\) in the
  actual image and equals \(F(t)\) at \(t\).
\end{itemize}

The same constant occurs in the support inequality and equality. The
hypothesis concerns a neighborhood in the image, so its pullback is a
neighborhood at any point of a curve mapping there. It does not assert
that an arbitrary CLD positive current has these data. At a zero
coordinate, a formula involving its logarithm must first be eliminated
or replaced by a finite floor as described below; cancellation of
undefined infinities is forbidden.

An \emph{actual rational AFS} function on \(Y=\operatorname{Spec}B\) is
\[u=\max\bigl(\lambda,c_i\log|f_i|+C_i:1\le i\le r\bigr),\tag{2.1}\]
where \(\lambda,C_i,c_i\) are rational, \(c_i>0\), and the list \(f_i\)
includes algebra generators of the \textbf{original} affine ring. Terms
with \(f_i=0\) have value \(-\infty\), while the floor is finite. Upper
growth means \(u\le H+C\) for an actual AFS exhaustion \(H\). It is not
the printed lower-growth convention.

For a finite-floor maximum, zeros are handled before tropicalizing. On a
compact neighborhood where the floor and all other weights are bounded
below by a finite number \(L\), a branch \(c\log|f|+C\) is inactive when
\(\log|f|\le r\), with rational \(r\) chosen so \(cr+C<L\). Choose a
second rational \(r'<r\). The domains \(|f|\le e^r\) and
\(|f|\ge e^{r'}\) cover, with a whole overlap. On the first domain the
branch can be removed. On the second \(f\) is invertible, with analytic
inverse, and its log is finite. Iterating over finitely many branches
gives finitely many charts. The same construction applies to finitely
many maxima and positive rational sums. A general selected weight is
covered by the unit-chart hypothesis itself; we make no unsupported
zero-coordinate reduction for arbitrary expressions such as
\(\log|z|-\log|w|\) at \((0,0)\).

There is also an explicit original ample reference for (2.1). Clear
slopes by an integer \(D\) so \(Dc_i\) are positive integers. Add, if
necessary, each original algebra generator \(z\) with coefficient
\(1/D\) and a sufficiently low rational intercept. It is dominated
globally by its existing branch and the floor: for \(0<\eta\le c\) and
suitably large \(M\),
\[\eta\log|z|-M\le\max(\lambda,c\log|z|+C).\tag{2.2}\] For \(\eta=c\)
this is immediate after reducing the intercept; for \(\eta<c\) the two
half-lines separated by the branch/floor intersection give a finite
maximum for the difference. This includes \(z=0\). The map given by
\(1\), the functions \(f_i^{Dc_i}\), and the original \(z\) is a closed
immersion of \(Y\) into an affine space because the \(z\) generate
\(B\). Its projective closure has \(Y\) as the actual standard affine
chart. The weighted FS metric on \(\mathcal O(1)/D\) has local weight
exactly \(u\). Its generic fiber is the original \(Y\), including
nonnormality. Merely using powers of generators would not prove this
chart assertion.

\subsection{3. Direct curve restriction and the order of
limits}\label{direct-curve-restriction-and-the-order-of-limits}

\textbf{Lemma 3.1.} A unit-chart rational contact weight pulls back to a
subharmonic function under every analytic map from a smooth strictly
analytic curve over a complete algebraically closed nontrivially valued
extension field.

Analytic base change pulls each unit and its inverse to a unit.
Evaluation of the pulled function equals evaluation at the projected
point, so the support inequality remains valid on the actual image. At a
curve point \(x\), the pulled support is
\[h_x=\sum_j a_j\log|f_j\circ\varphi|+C,\] which is harmonic by
Thuillier 2.3.21 (S6 LF3587--3624). It satisfies \(b\ge h_x\) nearby and
equality at \(x\).

To prove subharmonicity without a current pullback, let \(h\) be any
harmonic function and suppose \(b-h\) has a local maximum \(M\) at
\(x\). On a sufficiently small neighborhood,
\[h_x-h\le b-h\le M,\qquad (h_x-h)(x)=M.\] The harmonic maximum
principle, Thuillier 3.1.1 (S6 LF3730--3741), makes \(h_x-h\) constant
\(M\) on a smaller neighborhood. The two inequalities then force
\(b-h=M\) there. Thus \(b-h\) satisfies the strong local maximum
principle for every harmonic \(h\). Thuillier 3.1.12 (S6 LF3885--3896)
says exactly that a finite continuous function with this property is
subharmonic. This proves the lemma on arbitrary local analytic curves,
without extending their maps globally or replacing analytic units by
algebraic functions.

For a finite-floor maximum of positive logarithms, Thuillier 3.1.6 and
3.1.8 give the same conclusion, including zeros. Alternatively the
finite-chart construction of Section 2 reduces it to Lemma 3.1. Frame
changes add harmonic unit logarithms.

The order for a local decreasing sequence is precise. First pull back
every finite continuous contact approximant, on the inverse image of one
common neighborhood. Each is subharmonic by the lemma. Then take the
pointwise decreasing limit. Thuillier 3.1.8(ii), S6 LF3810--3818, gives
on each connected component either a subharmonic limit or the
identically \(-\infty\) function. A local lower bound excludes that
alternative for bounded inputs. No uniform convergence to an upper
semicontinuous discontinuous limit is asserted. In a later radial
argument an identically \(-\infty\) component is treated directly.

A useful but conditional variant is worth recording carefully. If actual
varying-class metric approximants restrict to weights
\(u_i+\epsilon_i h\), with \(u_i\downarrow u\),
\(\epsilon_i\downarrow0\), and a fixed continuous local reference \(h\),
shift \(h\) by a fixed constant to make it nonnegative on a relatively
compact domain. These weights then decrease to \(u\). This proves
restriction once that \textbf{actual reference identification} is
supplied. Numerical convergence of classes alone does not supply it.

\subsection{4. The actual strict cover, frames, and original
cocycle}\label{the-actual-strict-cover-frames-and-original-cocycle}

\textbf{Theorem 4.1.} Let \(k\) be algebraically closed and trivially
valued, \(X\) a reduced irreducible projective \(k\)-variety, and \(L\)
an ample rational line bundle. Suppose an actual globally glued
continuous metric on \(L\) has unit-chart rational contact weights,
allowing the finite-floor charts of Section 2. Then this metric is a
genuine continuous \(L\)-psh metric, obtained by uniform approximation
by genuine FS metrics on the same \(L\). Smoothness, countability, and
geometric unibranchness are not required in this partial theorem.

\textbf{Strict neighborhoods.} On an original affine frame chart, a
basic neighborhood of a seminorm is specified by finitely many upper and
lower bounds on regular functions. Choose finite rational bounds
strictly around the finite nonzero evaluations. At a zero choose only a
finite positive upper bound. Invert only functions for which a positive
lower bound is imposed. Add finite upper bounds on original affine
generators and on these inverses. The weak-inequality bounds form a
compact analytic domain with the tested point in its interior and
contained in the specified neighborhood. It is a domain in a
finite-radius closed affine polydisc with Laurent constraints. Finitely
many such compact neighborhoods cover \(X^{\rm an}\); shrink them into
the specified frame and contact charts. Unit logarithms are bounded
above and below on each.

Choose \(K=k((\varpi))\), \(|\varpi|=e^{-1/Q}\), with \(Q\) divisible by
the finitely many rational threshold denominators. This is a complete
nontrivially valued DVR field with residue field \(k\). The chosen
polydiscs and Laurent constraints become \textbf{strict} \(K\)-analytic
domains after base change. We do not invoke a strict-space theorem over
the original trivial field.

\textbf{Finite branch domains.} On one such compact strict domain the
unit logarithmic image is a finite rational polytope complex by
published Ducros Theorem 3.2, S9 LF2077--2155. Its proof first reduces
to algebraic rational domains by approximation that leaves norm images
unchanged, then uses the valued-field image theorem; its dimension
conclusion is only \(\le\dim X\). Refine the finite complex by all
equality hyperplanes of the selected affine expressions. On each
relative open convex cell, different expression classes have distinct
values everywhere. Continuity of the selection makes its expression
class locally constant; connectedness of the cell makes it constant.
Continuity extends that affine value to the closed cell.

Every closed-cell preimage is defined by finitely many rational affine
inequalities in unit logs. The finitely many cell constants are rational
multiples of the value-group generator. After one finite ramified
extension \(K'=k((\tau))\), \(\varpi=\tau^e\), and clearing integer
denominators, each inequality is of the form
\[|\tau^a\prod_j f_j^{m_j}|\le1,\qquad a,m_j\in\mathbb Z.\tag{4.1}\]
These are actual strict Laurent domains; equalities use both
inequalities. All negative powers are powers of units on the chosen
domain. There is a single common tensor denominator \(N\) for the
finitely many branches and charts. On every member \(V_\alpha\),
\[Nb_i=\log|a_\alpha|,\qquad a_\alpha\in\mathcal O(V_\alpha)^\times.\tag{4.2}\]
The original finite compact neighborhood cover is a G-cover because its
interiors cover. Each closed-cell subdivision is also a G-cover: for a
point, the finitely many members containing it have union a neighborhood
relative to the original domain, since the union of all other closed
members is closed and misses that point. This is the finite quasi-net
criterion, not a claim that every individual closed cell is a
neighborhood. The refinements are finite strict analytic domain covers
with actual analytic units.

\textbf{Frames and cocycle.} Fix the original algebraic frames \(e_i\)
and original transition functions \(e_i=g_{ij}e_j\) of \(L\). On
\(V_\alpha\) in chart \(i\) use \[s_\alpha=a_\alpha e_i^{\otimes N}.\]
Equation (4.2) gives \(\|s_\alpha\|=1\). On an overlap in charts
\(i,j\), the transition is exactly
\[s_\alpha/s_\beta=(a_\alpha/a_\beta)g_{ij}^{N}.\tag{4.3}\] It is an
invertible analytic function of norm one, and satisfies the cocycle
identity because it is obtained from the original algebraic cocycle by a
specified change of frames. It defines the original \(L^{\otimes N}\)
with its fixed identification, not an unrelated bundle inferred from a
tropical plot.

GM Proposition 2.10, S4 LF347--378, now applies. Its proof realizes the
genuine G-cover as formal opens, adjoins the norm-one transition
functions and their inverses, and uses one dominating admissible blowup
to make them formal units. Gluing by (4.3) produces a formal model of
the \textbf{same} bundle and metric. All prerequisites have been
supplied before this application.

\textbf{Nefness on all pieces.} Base change to
\(\Omega=\widehat{\overline{K'}}\). On every smooth analytic curve map
into any original weight chart, Section 3 proves subharmonicity of the
actual continuous weight. GR Proposition 5.1 and Corollary 5.4, S8
LF818--923, therefore give semipositivity of the already represented PL
weight. The criterion allows local analytic curve maps. Its reduction to
closed local curves uses normalization, and its residue-curve test uses
GR Theorem 4.1 with a good compact strict analytic space, a formal
model, and a point reducing to the generic point of the required residue
curve. That theorem produces a curve in a neighborhood of a possibly
moved point with the same reduction; it does not extend the curve
globally. Its statement and proof are retained at S8 LF489--800.

Every point on a gluing overlap and collar is tested, since the global
metric has the stipulated contact weights there as well. For a maximum
of two contact weights use the support of a winning weight; on a whole
reference-winning collar the metric equals the reference. Equality
merely on the boundary would not establish this antecedent. GM Lemma 3.3
detects nefness back on the already represented \(K'\)-model. GM
Proposition 3.11 then upgrades local nefness to global nefness. Its
proof tests each proper residual curve at a point reducing to its
generic point, uses the interior condition and a common dominating
model, and applies the projection formula (S4 LF754--782). No
unrestricted CLD current pullback or converse is used. GR Remark 1.8
expressly says that converse is unknown (S8 LF162--164).

\textbf{Algebraization and same-class approximation.} GM Remark 2.6 and
Lemma 2.4 algebraize on a projective model dominating the original
projective product model. The latter proof algebraizes a vertical formal
blowup ideal by formal GAGA. Its generic fiber remains \(X_{K'}\), not
its absolute normalization. Nefness is unchanged on domination and
descends by the projection formula.

BE Theorem 7.8, S7 LF2334--2352, applies over the nontrivially valued
\(K'\) to the ample \(L_{K'}\). On a dominating projective model choose
an ample model \(\mathcal A\) of the same generic bundle. Then
\[\mathcal L_\epsilon=(1-\epsilon)\mathcal L+\epsilon\mathcal A\] is
ample for positive rational \(\epsilon\), has generic bundle exactly
\(L_{K'}\), and its metric tends uniformly to the given metric.
Semiample model metrics are pure FS metrics. This states the actual
argument using one consistent ample bundle; the printed BE proof
switches its letter from \(A\) to \(H\). It supplies same-class FS
approximation, not singular EP.

\subsubsection{4.2. Exact Gauss coefficient cutoff and
descent}\label{exact-gauss-coefficient-cutoff-and-descent}

Give \(K'\) its rational step \(h=-\log|\tau|>0\). For every
\(x\in X^{\rm an}\), use the independent-variable Gauss completion of
\(\mathcal H(x)(T)\) with \(|T|=e^{-h}\) and embed \(\tau\mapsto T\).
The induced point \(\sigma(x)\in X_{K'}^{\rm an}\) restricts to \(x\).
For Laurent sections expanded in a fixed finite \(k\)-basis,
\[\left|\sum_{r\ge r_0}\tau^r f_r\right|_{\sigma(x)}
 =\max_{r\ge r_0}e^{-rh}|f_r(x)|.\tag{4.4}\] Independence prevents
cancellation even when the chosen radius already lies in the value group
of \(\mathcal H(x)\). The tails tend uniformly to zero on compact affine
charts, proving continuity and consistency on overlaps. This is the
Gauss section of BJ Remark 1.9 (S1 LF724--757), with a rescaled rational
value-group step.

For a pure FS generating system \(s_j\in H^0(X_{K'},mL_{K'})\), flat
field base change of \(H^0\) gives Laurent expansions
\(s_j=\sum_{r\ge r_0}\tau^r f_{jr}\) with \(f_{jr}\in H^0(X,mL)\). Let
\(h_0\) be the canonical trivial metric on \(L\). Every \(k\)-section
has norm at most one for \(h_0\): at the valuation-ring center of \(x\)
it is regular in a center frame, and its coefficient lies in that
valuation ring. Thus the pulled FS potential is
\[\sigma^*\psi=\frac1m\max_{j,r}
 \bigl(\log\|f_{jr}\|_{mh_0}-rh\bigr).\tag{4.5}\] The left side is
finite continuous on compact \(X^{\rm an}\) and has a finite lower bound
\(b\). Every term with \(-rh/m<b\) is uniformly below this lower bound,
so all sufficiently large \(r\) may be removed \textbf{exactly}. There
are finitely many remaining coefficients. They generate: at a common
algebraic zero their terms are all \(-\infty\) and the discarded tail is
strictly below \(b\), contradicting (4.5). All weights \(-rh\) are
rational. This produces a genuine finite FS metric on the same original
\(L\); no approximate coefficient cutoff or generic-fiber normalization
is involved.

Uniform errors decrease under pullback by \(\sigma\). Hence the original
continuous metric is a uniform limit of these genuine FS metrics. BE
Corollary 7.6, or BJ ample-class decreasing closure, gives genuine
membership. Explicitly choose uniform errors \(\epsilon_j=4^{-j}\) and
shifts \[B_j=2\sum_{i\ge j}\epsilon_i.\] If \(F_j\) has error
\(\le\epsilon_j\) relative to the target, then
\[F_j+B_j-(F_{j+1}+B_{j+1})\ge\epsilon_j-\epsilon_{j+1}>0.\] Thus a
same-class decreasing FS sequence is also available for this continuous
target. This completes Theorem 4.1.

GM 2.7 alone extends a represented metric from a compact domain, with no
preservation of nefness. Theorem 4.1 uses global gluing and all-point
curve tests instead. It does not silently replace the needed compact PL
extension or bounded product independence in Gate B.

\subsection{5. Canonical determinant weights without
flatness}\label{canonical-determinant-weights-without-flatness}

\textbf{Theorem 5.1.} Let \(k\) be a complete non-Archimedean field of
characteristic zero, \(A\) a normal integral finitely generated
\(k\)-algebra, and \(B\) an integral finite \(A\)-algebra with
\(A\hookrightarrow B\). Put \(F=\operatorname{Frac}A\),
\(E=\operatorname{Frac}B\), and \(d=[E:F]\). For every
\(s\in(\operatorname{Spec}A)^{\rm an}\) there are canonical positive
integers
\[m_y(s),\qquad y\in\pi^{-1}(s),\qquad \sum_y m_y(s)=d,\tag{5.1}\]
determined by the generic determinant norm. They agree with length times
residue degree where \(B\) is locally free. They need not equal the
dimensions of the nonflat special-fiber algebra factors. Every actual
rational AFS input has an actual AFS trace on the whole original base.

\subsubsection{5.1. Extending the polynomial law, including all
relations}\label{extending-the-polynomial-law-including-all-relations}

For \(b_1,\ldots,b_r\in B\), define
\[P(T_1,\ldots,T_r)=\det_F\bigl(\operatorname{mult}_{\sum T_i b_i}:E\otimes_F F[T]\to E\otimes_F F[T]\bigr).\tag{5.2}\]
The universal element is integral over \(A[T]\). Over an algebraic
closure of \(F(T)\) its eigenvalues are the field embeddings of that
integral element. Their elementary symmetric functions, including every
characteristic-polynomial coefficient, are integral over \(A[T]\) and
lie in \(F[T]\subset F(T)\). Normality of \(A[T]\) puts them in
\(A[T]\). This extends (5.2) across every nonflat base point.

This is a well-defined norm law on special fibers, not just a formula
for a chosen tuple. Choose finite \(A\)-module generators of \(B\) and a
finite presentation \(A^r/R\simeq B\). For any relation vector in \(R\),
the determinant is invariant under translating its variables by that
vector. The identities hold generically, hence coefficientwise in \(A\).
After tensoring to a field \(H\), the kernel of \(H^r\to B\otimes_AH\)
is the image of \(R\otimes_AH\), by right exactness. The translated
identities therefore imply well-definedness for every new fiber
relation. Multiplicativity, homogeneity of degree \(d\), the value
\(N(1)=1\), and the full Cayley--Hamilton identity also specialize,
since they are polynomial identities on finitely many universal module
generators. This remains valid if the fiber module has dimension larger
than \(d\).

\subsubsection{5.2. Factorization and strict positivity at every fiber
point}\label{factorization-and-strict-positivity-at-every-fiber-point}

Put \(H=\mathcal H(s)\) and \(C=B\otimes_AH\). Over an algebraic closure
\(\bar H\), write the finite algebra as a product of Artinian local
algebras, with residue characters
\(\chi_1,\ldots,\chi_q:C\otimes_H\bar H\to\bar H\) and nilradical \(J\).
On the unit group the specialized norm is a regular algebraic character.
In characteristic zero, the finite polynomial logarithm gives an
isomorphism \(1+J\simeq(J,+)\) because the algebra is commutative and
nilpotent. An additive affine space has no nontrivial regular character
into \(\mathbb G_m\): a character and its inverse are polynomial
functions whose product is one, so both are constant. On the residue
torus \((\mathbb G_m)^q\), a regular character is a Laurent monomial.
Polynomial extension to the whole algebra makes its exponents
nonnegative. Homogeneity gives
\[N_s(c)=\prod_{j=1}^q\chi_j(c)^{n_j},\qquad n_j\ge0,\qquad\sum_j n_j=d.\tag{5.3}\]
For completeness, the equality holds first on the dense unit open and
hence as a polynomial on the whole underlying vector space. The
torus-character assertion follows by comparing coefficients in
\(p(xy)=p(x)p(y)\) in the Laurent polynomial ring: a nonzero group-like
element has a single monomial with coefficient one.

Every \(n_j\) is \textbf{strictly positive}. If \(n_j=0\), take the
idempotent \(c\) with residue zero in that factor and residue one in
every other factor. Equation (5.3) would give \(N_s(c)=1\). The
specialized degree-\(d\) characteristic polynomial annihilates \(c\) and
has constant term \((-1)^dN_s(c)\ne0\). Solving that identity for an
inverse makes \(c\) invertible, which contradicts its zero residue. This
also proves why an invisible point cannot be assigned weight zero.

Galois invariance makes \(n_j\) equal on the embeddings of each residue
field \(H_y\) of a local factor of \(C\). Characteristic zero makes
these finite residue extensions separable. A finite extension of a
complete non-Archimedean field has a unique extending absolute value.
Set \[m_y(s)=n_j[H_y:H].\tag{5.4}\] These are the positive integers of
(5.1). They are intrinsic and input-independent. If \(B\) is locally
free, (5.2) is its ordinary determinant law; an Artinian composition
series gives \(n_j\) equal to the local Artinian length. There is no
local-free hypothesis elsewhere in the argument.

Under any complete field extension of \(H\), the residue factors split,
with the sum of their new residue degrees equal to the original degree,
and their exponents are unchanged. Thus the refined weights aggregate to
(5.4). This proves base-change compatibility of weighted sums of any
pulled-back function, also for value \(-\infty\) when the functions are
locally bounded above. It uses the actual specialized polynomial law,
rather than an affine regularization theorem over \(H\).

\subsubsection{5.3. Exact AFS formula on the whole
base}\label{exact-afs-formula-on-the-whole-base}

For (2.1), choose \(D\) clearing the \(c_i\), put
\[b_0=1,\quad b_i=f_i^{Dc_i},\quad w_0=D\lambda,\quad w_i=DC_i,\] and
expand the actual polynomial
\[P(T)=N\bigl(\sum_{i=0}^rT_i b_i\bigr)=\sum_{|\beta|=d}a_\beta T^\beta\in A[T].\]
Give \(T_i\) weighted Gauss radius \(e^{w_i}\). The Gauss norm is
multiplicative at arbitrary positive real radii. One elementary proof
chooses the lexicographically largest maximal-weight monomial in each
polynomial: their sum exponent has a unique maximal contribution among
maximal-weight terms, and no lower-weight term can cancel it. Apply this
after complete algebraic closure and (5.3). Conjugate residue values
have equal absolute value. Taking logarithms gives the exact identity
\[N_\pi u(s):=\sum_y m_y(s)u(y)
 =\max_{|\beta|=d}\left(\frac1D\log|a_\beta(s)|+\frac1D\sum_i\beta_iw_i\right).\tag{5.5}\]
The coefficient of \(T_0^d\) is one, providing the finite floor
\(d\lambda\). Zero coefficients are discarded. All other branches use
regular functions in the original \(A\), so the displayed maximum is
continuous everywhere, including the nonflat locus.

To retain any fixed original base generator \(g\), express \(g\) as a
polynomial in the source generators in (2.1). Degree bounds, the
positive source slopes, and the finite floor give rational \(\eta>0\),
\(M\) with \[u\ge\eta\log|g\circ\pi|-M.\] Since the weights are positive
and total \(d\), append the dominated branch \(d\eta\log|g|-dM\) to
(5.5). The value does not change. Doing this for all base generators
makes the result an actual AFS, with the original ring and generators.
Section 2 supplies its actual ample-reference compactification. Thus
``genuine'' here has finite algebraic FS data, rather than a tropical
inverse-measure meaning.

\subsubsection{5.4. A nonflat example that separates weight from fiber
length}\label{a-nonflat-example-that-separates-weight-from-fiber-length}

Let \(k\) be any characteristic-zero field, \(A=k[x,y]\), \(I=(x,y)\),
and inside \(A[z]/(z^2-x)\) let \[B=A\oplus Iz.\] Multiplication
preserves this finite \(A\)-module: \((uz)(vz)=xuv\in A\). The ambient
quadratic ring is integral, so \(B\) is integral and has generic degree
two. At the origin, \[B/(x,y)B\simeq k\oplus (I/(x,y)I)z,\] which has
dimension three, with a two-dimensional square-zero ideal. Hence the
finite morphism is not flat there. Nevertheless \[N(a+uz)=a^2-xu^2,\] so
the specialized norm at the origin is \(a(0)^2\). There is one fiber
point of canonical weight two, not three.

For original source generators \(x,y,xz,yz\), the universal determinant
is \[P=(T_0+xT_1+yT_2)^2-x(xT_3+yT_4)^2.\tag{5.6}\] At the origin it is
\(T_0^2\), making an AFS floor \(\lambda\) trace to \(2\lambda\).
Multiplying the special-fiber algebra dimension by that floor would give
the wrong \(3\lambda\). This is an exact symbolic example, not a
numerical sample.

\subsection{6. Whole-base continuity, local integer sums, and finite
pullback}\label{whole-base-continuity-local-integer-sums-and-finite-pullback}

\textbf{Theorem 6.1.} With the canonical weights of Theorem 5.1, the
trace of any continuous real-valued function on \(Y^{\rm an}\) is
continuous on \(S^{\rm an}\). On any compact base set \(E\),
\[\|N_\pi a-N_\pi b\|_E\le d\|a-b\|_{\pi^{-1}E}.\tag{6.1}\] If
\(u_j\downarrow u\) are actual AFS and \(u\) is continuous, then their
traces converge locally uniformly, and each trace is an actual AFS on
the whole base.

The estimate follows pointwise from positivity and (5.1). A finite
morphism is proper, so the inverse image of a compact base neighborhood
is compact. Dini's theorem on that compact set applies when the limit is
continuous; (6.1) then proves locally uniform convergence of the genuine
AFS traces (5.5). This is the required whole-base continuity mechanism
for such inputs. It does not infer continuity of a discontinuous
bounded-decreasing limit.

For the stronger assertion concerning every continuous input,
differences of actual AFS form a rational vector lattice. Positive
rational sums and maxima are AFS after distributing the finite maxima
and clearing powers. Original generator branches can be appended by
(2.2). The identity \[\max(a-b,c-d)=\max(a+d,c+b)-(b+d)\] proves lattice
closure for differences. On a compact source set, this lattice contains
constants in restriction, by moving generator intercepts low enough
below a desired floor; it separates seminorms, by truncated log norms of
a distinguishing regular function. Arbitrary real scalars can be
approximated by rationals on that compact set.

Here is the needed lattice-density proof. For fixed \(x,y\) and
continuous \(f\), choose a lattice function with prescribed values at
\(x,y\) within a chosen error, using separation, scaling, and constants.
For fixed \(x\), finitely many such functions, indexed by \(y\), have
neighborhoods where they exceed \(f-\epsilon\) and these neighborhoods
cover the compact set. Their maximum exceeds \(f-\epsilon\) everywhere
and remains below \(f+\epsilon\) near \(x\). A finite collection of
these neighborhoods covers as \(x\) varies. The minimum of the
corresponding maxima lies between \(f-\epsilon\) and \(f+\epsilon\)
everywhere. It is still in the lattice. Thus differences of AFS
uniformly approximate \(f\) on the compact set. Their traces are
differences of continuous determinant formulas, and (6.1) gives
continuity of the limit. This assertion is continuity, not a positivity
theorem for arbitrary continuous inputs.

If a continuous input equals one on one clopen fiber piece and zero on
the complementary piece over an open base set, its trace is the sum of
the canonical positive integer weights in the first piece. Continuity
makes this integer sum \textbf{locally} constant. It can take different
values on different components of that open. This is precisely the
conclusion needed for a later shrink around one tested fiber.

Finite pullback also retains original source generators. If \(v_j\) is
AFS on \(S\), its pullback has base branches. Every source generator
\(z\) satisfies a monic integral equation \[z^m+\sum_{i<m}b_i z^i=0,\]
whence \[\log|z|\le\max_{i<m}\frac{\log|b_i|}{m-i}.\tag{6.2}\] For
\(z=0\) the statement is immediate. Otherwise, if its norm exceeded
every right-hand branch, the leading monomial would have strictly larger
norm than every other term, contradicting the equation. Base-generator
degree bounds and the positive branches and floor of each \(v_j\) let
one append a sufficiently small positive \(z\)-branch with a
sufficiently low rational intercept, dominated everywhere. Thus every
pulled function is actual AFS on the original \(B\), and the decreasing
limit is unchanged. The denominator \(m-i\) and the source generators
are indispensable.

\subsection{7. Radial identification conditional only on open-trace
extension}\label{radial-identification-conditional-only-on-open-trace-extension}

We now retain \(k\) algebraically closed, trivially valued, and
characteristic zero, \(S\) an affine open in \(\mathbb A_k^n\), and
\(Y\) integral affine finite surjective over \(S\). Say \(u\) has a
\emph{contact witness} if it is locally bounded above, is finite at the
generic trivial seminorm, and has one proper birational algebraic
morphism \(\mu:Z\to Y\) such that \(u\circ\mu^{\rm an}\) is locally a
decreasing sequence of the contact weights of Section 2 on common
neighborhoods.

\textbf{Theorem 7.1 (conditional extension identification).} Suppose
\(v\) is an actual AFS-regularizable function on all of \(S^{\rm an}\)
and, on some dense principal open \(W\subset S\) over which \(\pi\) is
étale and the witness is an isomorphism, \(v=N_\pi u\). Then
\[v(s)=\sum_{y\in\pi^{-1}(s)}m_y(s)u(y)\quad\text{for every }s\in S^{\rm an}.\tag{7.1}\]
The hypothesis that such a \(v\) exists is retained. It is the
open-trace regularization/removal output of Gate B, and is not proved by
the determinant formula. The theorem removes the separate residue-field
affine-regularization premise from the identification step.

\subsubsection{7.1. Curve slices agree with the canonical
weights}\label{curve-slices-agree-with-the-canonical-weights}

Fix \(s\) and set \(H=\mathcal H(s)\). The canonical \(H\)-rational lift
\(s'\) on \(S_H\) has finite fiber \(\operatorname{Spec}(B\otimes_AH)\),
whose analytic points are exactly the original fiber. Choose a line
through \(s'\) generically in \(W_H\). This exists because \(H\) is
infinite: for each polynomial defining the complement, its restriction
to a line in a variable direction is not identically zero for a
direction outside a proper algebraic subset. Intersect with \(S_H\) to
obtain an affine open of that line.

The scheme-theoretic preimage has pure dimension one. Indeed \(Y_H\) is
an integral \(n\)-dimensional finite-type scheme (geometric integrality
follows from algebraic closedness of \(k\)); the \(n-1\) line equations
give height at most \(n-1\) at its components, and finiteness gives
dimension at most one. Thus each reduced component dominates the line.
Normalize the reduced preimage. The resulting disjoint smooth curves
\(C\) over the perfect \(H\) map finitely and flatly to the line: their
coordinate modules are finite torsion-free over the local PIDs of the
smooth base curve. Over its generic point in \(W_H\) the map is the
original étale map, so total rank is \(d\).

The generic determinant norm of the source equals the determinant norm
on \(C\) for every original regular source function, because they agree
over that étale open. Both sets of coefficients are regular on the
normal line, so the equality holds on the whole line, including the
center. Hence (5.3), or uniqueness of its exponents, shows that the
length/residue weights of the normalized slice, aggregated over each
original fiber point, equal exactly the canonical weights \(m_y(s)\).
These weights were fixed before the slice or weak input was chosen.
Different admissible slices give the same result. No special-fiber
dimension from the nonflat source is used.

The generic lift of every component of \(C\) through the single witness
is defined over the isomorphism open. Properness of \(\mu_H\) extends it
over the DVR at each missing curve point, and separatedness gives
compatible uniqueness. It therefore gives an algebraic map \(C\to Z_H\).
This is the required local lifting criterion, with its hypotheses
checked; arbitrary local analytic curves are not asserted to extend
globally.

\subsubsection{7.2. Direct Gauss limits over the completed residue
field}\label{direct-gauss-limits-over-the-completed-residue-field}

If \(H\) is trivially valued, first extend to \(H((T))\), with its
nontrivial \(T\)-adic norm. Then in either case pass to a completed
algebraic closure \(\Omega\). Every original fiber value is preserved
under its rational lifts. Pullback of each supplied local contact
approximant through the algebraic lift above satisfies Lemma 3.1. Taking
the decreasing limit gives a subharmonic function, or the identically
\(-\infty\) function, on each source curve neighborhood. The actual AFS
sequence supplied for \(v\) restricts on the base line to finite
positive maxima of logs of regular functions and a finite floor, hence
gives the same alternative on a base neighborhood. This does not use an
AFS approximation theorem over \(H\) or \(\Omega\).

For a subharmonic function \(h\) on a disk about a rational point \(0\),
harmonic comparison on every sufficiently small closed disk gives
\[h(0)\le h(\xi_R),\] where \(\xi_R\) is its boundary Gauss point. The
Gauss points tend to \(0\), so upper semicontinuity gives the opposite
limsup inequality. Thus \[\lim_{R\downarrow0}h(\xi_R)=h(0),\tag{7.2}\]
also if \(h(0)=-\infty\). For an identically \(-\infty\) neighborhood
this is immediate.

At a point \(x\) of the smooth normalized source curve over the center,
choose an analytic coordinate \(z\). The finite map has the form
\[t=c z^e U(z),\qquad U(0)=1,\ c\ne0.\] Here \(e\) is a positive
integer. Residue characteristic zero ensures that after shrinking, the
binomial series for \(U^{1/e}\) converges; the coordinate
\(zU(z)^{1/e}\) has an invertible linear coefficient. The map is
therefore exactly \(t=cy^e\) on a smaller disk. Finite properness puts
every preimage of a sufficiently small target disk in these finitely
many source disks.

Choose target radii \(R\in|\Omega^\times|\) tending to zero; then
\((R/|c|)^{1/e}\in|\Omega^\times|\). Over the target Gauss point there
is one source Gauss point in each such disk, of weight \(e\). To check
the degree, scale the coordinates to radius one: the residual polynomial
is \(Y^e-\alpha T\) over \(\widetilde\Omega(T)\), irreducible by
Eisenstein at \(T\). A monic reduction argument gives degree \(e\) in
the completed residue field. The known degree of the disk map and the
unique norm extension exclude another point. This is an explicit
Gauss-fiber computation, with no punctured-limit assertion.

The complement of \(W_H\) on the line is finite. Small noncentral Gauss
points therefore lie in \(W_H\). At each such point the pulled open
identity gives
\[v(\xi_R)=\sum_x e_x (u\circ p)(\xi_{(R/|c_x|)^{1/e_x}}).\tag{7.3}\]
Pass to the limit using (7.2). There are finitely many summands. If one
tends to \(-\infty\), all others are bounded above on these
neighborhoods, so the sum tends to \(-\infty\) without subtracting
infinities. The central slice weights are the canonical weights by
Section 7.1, and residue-field refinements aggregate by Theorem 5.1.
Equation (7.1) follows.

Consequently, after Gate B supplies the actual extension \(v\), the
trace is identified on the \textbf{whole} smooth base and is compatible
with the single original witness. Combined with Section 6, it gives the
local integer sums, degree estimate, and actual finite pullback needed
by the conditional split-piece proof. It does not by itself supply that
extension or the singular envelope theorem.

\subsection{8. Bounded local products and the remaining smooth
gate}\label{bounded-local-products-and-the-remaining-smooth-gate}

Theorem 4.1 removes the representation and positivity circularity for
its explicit continuous contact metrics. The following construction
proves a bounded local product statement for that exact decreasing
class. It does not accept the embedded round70 chain by citation.

The literal CLD Proposition 8.7.2 assumes continuous functions
converging \textbf{locally uniformly}. Its proof multiplies by
continuous functions using Lemma 8.5.1 and inducts after applying
\(d'd''\) (S3 LF10432--10449). A bounded decreasing local sequence with
a discontinuous limit does not meet that premise. BE Theorem 8.4 is
likewise a uniform extension on continuous psh-regularizable metrics,
with exact fundamental-cycle mass in 8.5 (S7 LF2611--2779). The
construction below uses genuine global decreasing convergence instead.

\subsubsection{8.1. Common-collar bounded products and weighted
limits}\label{common-collar-bounded-products-and-weighted-limits}

\textbf{Proposition 8.1.} On an analytic open of the original projective
variety, let a finite tuple \(u_1,\ldots,u_n\) be locally bounded
decreasing limits of the contact weights of Section 2 on common
neighborhoods. There is an intrinsic positive mixed local measure,
agreeing with the continuous local product, obtained by limits of the
supplied continuous products. The limit is independent of their supplied
sequences, common collars and auxiliary ample references. The same
conclusion holds for a bounded local contact test \(u_0\) multiplying
the measure, with its supplied decreasing approximation. No local
comparison or IBP assertion is included.

We first construct actual collar data around any point inside the common
neighborhood. Fix a very ample line bundle \(H\) on \(X\) and a section
\(s\) whose affine chart contains the point's algebraic support. Basic
analytic neighborhoods there have finitely many regular constraints with
finite rational strict margins. Each inverted nonvanishing function is
\(\tau_i/s^{d_i}\), by ordinary section localization; this uses no
normality. The section \(a=s\prod_i\tau_i\) of a positive power of \(H\)
gives the exact principal chart where all those inverses are regular.
Each new constraint can be written \(F_j/a^m\), after increasing a
common \(m\). Increase it until that power of \(H\) is very ample and
append a complete section system with sufficiently loose bounds. Its
ratios by \(a^m\) generate the original affine ring, by the standard
chart of the closed projective embedding. A rational FS local weight
\[v=\frac1q\max\{q\lambda,\log|F_j/a^m|-C_j\}\] with finite rational
thresholds and negative floor \(\lambda\) has compact \(Q=\{v\le0\}\)
inside the specified neighborhood. Choose a rational \(\epsilon>0\)
below all divided strict margins and with \(\lambda<-\epsilon\). Then
the interior of \(P=\{v\le-\epsilon\}\) contains the point. Here \(q\)
is the common tensor exponent of \(H\); the thresholds are \(C_j\), not
\(qC_j\). This is an actual FS metric on \(H\), with global model
potential \(p\) relative to its canonical trivial metric. Write
\(\theta=\omega+dd^cp\) for its curvature, where \(\omega=c_1(H)\)
denotes that fixed original reference.

The section-localization assertion follows by ordinary module
localization on a finite affine trivializing cover and by its finite
equalizers for matching sections: filtered localization commutes with
those equalizers. Thus
\(\Gamma(D(s),\mathcal O_X)=\varinjlim_d H^0(X,H^d)/s^d\), retaining the
original nonnormal ring.

Choose \(M\) bounding all the local limits, including the test if
present, on a common neighborhood of \(Q\). Choose rational \(A\) with
\(A\epsilon>2M+2\). Then \(h=Av+M+1\) is below \(-M\) on \(P\) and at
least \(M+1/2\) on the compact strip \(Q\cap\{v\ge-1/(2A)\}\). For each
supplied decreasing continuous sequence, its closed superlevel sets
\(\{u_{i,j}\ge M+1/4\}\) on this strip decrease to the empty set.
Compactness gives one common eventual tail for all finitely many
sequences with \(u_{i,j}<M+1/4<h\) on the \textbf{whole} strip. This
invokes no Dini theorem for a discontinuous limit.

Set \(\Psi_{i,j}=\max(u_{i,j}-h,0)\) on \(Q\), zero outside. They agree
with zero on the whole strip, hence on a neighborhood of the gluing
boundary. Their metric weights are \(\max(u_{i,j}-M-1,Av)\) inside and
\(Av\) on the reference-winning strip. These have the contact weights
and whole collars required by Theorem 4.1. Their actual class is
\(A\theta\), on the same algebraic \(H\) and original cocycle after
clearing \(A\); no desired weak membership is used. Their explicit
bounded decreasing limits \(\Psi_i\) are genuine by decreasing closure.
On \(P\) put \[\Phi_i=\Psi_i/A+p,\qquad G=v-p+(M+1)/A,\] so
\(u_i/A=\Phi_i+G\). The actual global metric with local weight
\(v+\Psi_i/A\) equals \(u_i/A-(M+1)/A\) on \(P\). Its
canonical-reference potential is \(\Phi_i\) in the original ample class
\(\omega\), not \(\Psi_i/A\) alone. Every \(\Phi_i\) is bounded genuine
finite energy. Define the local mixed measure on the interior of \(P\)
by \[A^n M^{\rm raw}_\omega(\Phi_1,\ldots,\Phi_n).\] The factor \(A^n\)
accounts for scaling from \(u_i/A\); no extra factor is inserted if one
instead retains the unscaled global metrics in class \(A\theta\).

For continuous inputs this is the continuous local product. Theorem 4.1
supplies genuine same-bundle FS approximation; BE7.14 smooths the FS
metrics by smooth maxima (S7 LF2370--2379). BE8.18 identifies their
model Deligne intersections, with fundamental cycles, by induction using
a relatively regular section whose norm is one at every model Shilov
point (LF3125--3202). Along with uniform continuity, this is BJ
Remark7.23's genuine/CLD continuous compatibility (S1 LF4126--4129).
Continuous locality is S1 Remark7.42 (LF4546--4551). These applications
stay on the inner open, avoiding the boundaries of auxiliary affinoid
domains; the represented-nef-to-positive-current implication is also
stated and proved at S10 LF2547--2589.

Genuine decreasing continuity of the finite-energy mixed operator gives
convergence for the supplied global truncations, hence for every
compactly supported continuous test in the inner open. To prove
independence, compare one supplied sequence in its original collar and
in a third smaller shared collar about the tested point. At every finite
stage their continuous local products agree there: both have the same
actual local weights up to the stated constant and scale. Pass to both
genuine decreasing limits against compact tests. For two different
supplied sequences in that shared collar, their explicit global limits
are the same function \(\max(u_i-h,0)\) there, zero outside. The genuine
global operator gives the same measure. Comparing each sequence with the
shared collar proves independence of sequences and collars. Changes of
reference use the same continuous equality and shared-collar argument,
rather than comparing unequal numerical representatives. The identities
hold near every point and therefore glue the locally finite positive
measures on the analytic open.

For a bounded decreasing contact test use the same construction for
\(i=0\). On the inner open, \(u_{0,j}=A\Phi_{0,j}+AG\) with fixed
continuous \(G\) in the chosen algebraic frame. BJ7.21 applies to the
bounded genuine \(\Phi_{0,j}\) and the other bounded genuine factors;
the \(G\) term passes by ordinary weak convergence against compact
cutoffs. The sign in the scalar pairing identity is the pairing
containing \(\Phi_0\) \textbf{minus} the zero-test pairing, as printed
correctly at LF4108--4113; LF4119--4123 reverse it. With the correct
order both pairings converge. Translate the test below zero and apply
the positive weighted-measure Lemma7.22 to its negative to get weighted
weak convergence. This rederives the application without treating a
difference of two decreasing functions as decreasing. The same
shared-collar comparison establishes independence of the weighted limit.
This proves Proposition 8.1.

\subsubsection{8.2. What the bounded local product does not yet
establish}\label{what-the-bounded-local-product-does-not-yet-establish}

To complete Gate B one must still transfer this bounded result through
signed local PL extensions, construct compact PL extensions and
subordinate rational PL partitions with controlled support, and prove
compact IBP with support retained in a common collar. A telescoping sum
must then prove \textbf{raw} global bounded-local mixed mass, and
common-reference locality must give comparison and domination. These
steps precede smooth local/global equality; domination of an auxiliary
weak cone cannot be assumed from genuine global domination.

The attempted round70 smooth growth proof has a clear logical
prerequisite. Its potential \[\Psi_M=\max(\bar\pi^*u-cA\bar\pi^*g,M)\]
on a resolved projective witness has a constant boundary collar from the
\textbf{upper} bound \(u\le Ag\), for \(c>1\). To call this potential
genuinely psh in the nef pullback class, it uses smooth local/global
equality. That equality depends on the preceding bounded products, local
mass, and local domination. BJ Corollary 12.18 starts with
\(\phi\in\operatorname{PSH}(\theta)\), literally at S1 LF6259--6265; it
cannot establish this first membership. Once membership is established,
its same-ample-class sequence and original-generator conversion can be
used. The growth argument is therefore not an independent proof of the
missing bounded local theory.

Exact smooth principal-open removal is also unfinished. One must first
regularize the controlled \(v+c\log|f|\), then use smooth tail
envelopes, and then identify removed-set values on the single original
resolved witness through a valid common-reference divisorial-order
argument. Upper semicontinuity and equality on a dense open do not prove
those values. Without this reconstruction, the extension premise in
Theorem 7.1 remains indispensable.

This is a missing proof, not a disproof of the corrected UPPER-growth
theorem. The previously supplied quadratic counterexample addresses only
the printed LOWER-growth convention and is not used to reject singular
EP.

\subsection{9. The new envelope source and the actual global equation
theory}\label{the-new-envelope-source-and-the-actual-global-equation-theory}

Wu v1, S2 LF2697--2755, states EP for geometrically unibranch varieties
in its retained countable discretely/trivially valued
equicharacteristic-zero setting. Its proof explicitly invokes
Proposition 2.24 and local/global Theorem 6.5 at LF2705--2717, then a
countable-field reduction. The statements and entire proof are retained
literally. They do not eliminate the need to reconstruct Gate B, the
singular Runge induction, and its own finite-data extension. No
independent correctness verdict is inferred from the public v1 preprint.

BJ v4 defines reduced possibly nonnormal varieties at S1 LF435--469, and
Section 7 resumes general projective varieties at LF3649--3658. Its
actual finite-energy operator, comparison and uniqueness have this
scope. It does not use absolute normalization of a test configuration:
relative integral closure takes place inside its prescribed generic
fiber (LF759--849). Any later base-change application must preserve that
original nonnormal generic fiber and its vertical \(R_1,S_2\)
conditions.

For the equation retain \(n=\dim X\), \(\omega=c_1(L)\), \(V=(L^n)>0\),
and BJ's genuine global cone. Its normalized energy is
\[E_\omega(\phi)=\frac{\langle(\omega,\phi)^{n+1}\rangle}{(n+1)V},\qquad
\mathcal E^1=\{\phi\in\operatorname{PSH}(\omega):E_\omega(\phi)>-\infty\}.\]
The genuine raw model measure on a relatively integrally closed
determination is
\[M^{\rm raw}=\sum_E b_E\bigl((\omega_{\mathcal X}+D)|_E\bigr)^n\delta_{v_E},
\quad b_E=\operatorname{ord}_E(t),\quad v_E=b_E^{-1}\operatorname{ord}_E|_{k(X)}.\]
Its mass is \(V\), with the central fundamental cycle and \(b_E\)
retained. BJ Theorem 7.18 constructs the actual extended measure by
pairing with genuine PL tests and Radon representation; decreasing and
weighted convergence follow from the corresponding pairings (S1
LF3980--4148). We cite this primary theorem for its actual scope,
without importing round72's proposed corrections or proof as a new black
box.

A target is a positive Radon \(\nu\) of mass \(V\), with \(\mu=\nu/V\)
satisfying \[E_\omega^\vee(\mu)=\sup_{\psi\in\mathcal E^1}
\left(E_\omega(\psi)-\int\psi\,d\mu\right)<\infty.\] On irreducible
\(X\), BJ 9.20 equates this to integrability of every genuine
finite-energy potential (S1 LF5266--5301). There is no support
restriction. BJ Corollaries 10.4 and 10.6 give uniqueness modulo
constants and domination before EP (LF5440--5477).

The decisive source theorem is BJ \textbf{12.8}, not its smooth endpoint
12.9. The former states that EP is equivalent to onto finite-energy MA.
Its proof is retained at S1 LF6092--6141: choose a normalized dual
maximizing sequence; EP gives weak compactness, bounded energy gives a
finite-energy limit, 9.19 gives target integral convergence on that
bounded-energy set, and energy upper semicontinuity gives an actual
maximizer. Genuine differentiation and duality identify its MA. The raw
objective \(VE(\phi)-\int\phi\,d\nu\) is constant-invariant exactly at
mass \(V\). This source proof is conditional on EP for the singular
space. Constructing the operator or strict concavity alone does not
provide its compactness hypothesis.

\subsubsection{9.1. Exact outstanding dependencies after the new partial
results}\label{exact-outstanding-dependencies-after-the-new-partial-results}

The following remain unproved in this deliverable:

\begin{enumerate}
\def\labelenumi{\arabic{enumi}.}
\tightlist
\item
  Gate B in the exact form of Section 8, including compact PL extension
  with controlled support and original cocycle, signed compact IBP,
  global bounded-local mass/comparison/domination, smooth local/global
  equality, weak UPPER-growth regularization, controlled approximation,
  and removed-set equality on the original witness.
\item
  The extension premise in Theorem 7.1 for every weak upper-growth
  input. Theorem 5.1 fixes nonflat weights unconditionally, and Theorem
  7.1 identifies an existing regularizable extension, but neither
  constructs that extension without Gate B.
\item
  The full conditional split-piece, descending-\(Z\), Hartogs and
  reduced-singular-strata chain, now using the actual metric and
  canonical trace results. Runge must precede the strict coefficient
  radius formula, and the final empty-\(Z\) step and countable
  weak/genuine equality must be proved. The embedded round72
  implications remain unreviewed; we do not declare them automatically
  valid after improving two antecedents.
\item
  Actual varying-class/reference compatibility for resolved nef fiber
  constancy, geometrically connected fiber and curve-morphism
  applications, and proper divisorial minimum-usc descent. The proposed
  Stacks applications are supplied paraphrases without pinned pages; no
  independent access to them is claimed here. Theorem 3.1 gives
  restrictions for supplied contact approximants, but does not turn
  every BJ numerical nef approximant into the required
  original-reference metric without further argument.
\item
  Countable rational geometrically unibranch EP, followed by
  arbitrary-field finite-data descent. Only individual finite PL/model
  obstacles and competitors can descend to suitable countable fields.
  Positivity, actual relative generic-fiber closure, raw mass,
  orthogonality, and domination must be proved compatible with that
  extension before contact integrals compare all competitors. Arbitrary
  PSH competitors do not descend simultaneously.
\item
  Onto normalization only after target-side approaching EP, degree-one
  fundamental-cycle projection, and unbounded energy/MA compatibility;
  then actual singular variational compactness, differentiation and
  attainment for every target above. A homeomorphism of analytifications
  alone proves none of these onto analytic assertions.
\end{enumerate}

The trace result does not define an auxiliary inverse measure and is not
substituted for the genuine MA measure. The operator and conditional
equation description retains the original ample polarization and
trivially valued characteristic-zero field hypotheses. No arbitrary
complete-field equation theorem is asserted. The original completion
criterion remains unmet.

\subsection{10. Reproduction and
evidence}\label{reproduction-and-evidence}

Run \texttt{python3\ rerun.py} in this workspace. It verifies
byte-identical input preservation and all 23 supplied primary-file
hashes, recreates 50 literal units with inclusive one-based coordinates
defined by UTF-8 byte decode and
\texttt{split(\textquotesingle{}\textbackslash{}n\textquotesingle{})},
runs bounded exact diagnostics, and uses existing Pandoc and pdfLaTeX
for two passes. Both pass logs, stdout/stderr, commands, tool versions,
and output hashes are retained. No installation, other
worker/model/subagent, other campaign/private/root read, external
message, source-cache edit, runtime/auth change, or network access is
used.

The finite diagnostics cover explicit determinant polynomials, the
nonflat rank-two/fiber-length-three example, degree bounds for fixed
integer weight lists, retained-generator monic bounds, and finite
collar/shift inequalities. They are exhaustive only on their declared
finite grids. They establish no infinite analytic theorem. Successful
compilation, matching hashes and finite checks are observations,
distinct from mathematical validation.

\texttt{status.json} has exactly \texttt{id}, \texttt{status},
\texttt{exact\_result}, \texttt{proof\_file}, and
\texttt{concrete\_gaps}. \texttt{canonical-notes.json} contains four
backward-supported notes. The standalone artifacts are
\texttt{P16.1-paper.tex} and \texttt{P16.1-paper.pdf}. No resource-heavy
computation was required; no operator-provided supervised compute
service was exposed or used.

\subsection{Primary bindings}\label{primary-bindings}

S1: Boucksom--Jonsson, \emph{Global pluripotential theory over a
trivially valued field}, 1801.08229v4. S2: Lyuhui Wu, \emph{On the local
theory of non-archimedean psh functions}, 2609.27610v1. S3:
Chambert-Loir--Ducros, 1204.6277v2. S4: Gubler--Martin, \emph{On Zhang's
semipositive metrics}, 1608.08030v3. S5: Wanner, \emph{Comparison of two
notions of subharmonicity on non-Archimedean curves}, 1801.04713v3. S6:
Thuillier, thesis tel-00010990v1. S7: Boucksom--Eriksson, \emph{Spaces
of norms, determinant of cohomology and Fekete points in non-Archimedean
geometry}, 1805.01016v3. S8: Gubler--Rabinoff, \emph{On semipositive
piecewise linear functions in non-archimedean analytic geometry},
2503.05513v1. S9: Ducros, published \emph{Confluentes Mathematici} 4(4),
1250007, DOI10.1142/S1793744212500077. S10: Gubler--Jell--Rabinoff,
\emph{Forms on Berkovich spaces based on harmonic tropicalizations},
2111.05741v3.

Every exact supplied PDF/text/HTML hash and filename is in
\texttt{source-extracts/source-provenance.json}; authorized
byte-identical local copies are in \texttt{sources/}. Literal statements
and pertinent proof bodies, their versions and coordinates, and concrete
scope limits are retained. Retention is not a claim that every full
source body was independently audited. All proofs and source
applications remain unverified drafts.

\end{document}
